
% =========================================================
% CHAPTER 1: INTRODUCTION
% =========================================================
\chapter{Introduction}

\section{Introduction}
The rise of the digital economy in Pakistan has been driven by the widespread adoption of fintech platforms such as Easypaisa, SadaPay, and JazzCash~\cite{baig2022future}. These services have empowered millions of users by simplifying payments and providing instant access to their complete transaction history. This readily available data holds the potential to offer valuable insights into personal spending habits. However, for most users, this potential remains untapped as the data is typically provided in raw formats (e.g., CSV files) without sophisticated analytical tools.

This report documents the Smart Transaction Sorter and Analytics system, a web-based application designed to unlock the value hidden within this raw financial data. The project's academic objective is to develop a functional, standalone tool that serves as a research-based feasibility study. The resulting application is structured as a practical blueprint for an intelligent analytics module. The long-term vision, while outside the immediate scope of this project, is that such a module could be pitched for integration into the very fintech platforms it complements, thereby closing the significant analytical gap in the current market.

\section{Problem Statement}
The problem is the absence of an integrated Business Intelligence (BI) and analytics layer in Pakistan digital wallets, prompting many users to request spending breakdown features online~\cite{reddit2023sadapay}. People export their data, move it into another tool, and tag each transaction by hand. It is a manual and time-consuming process.

\section{Proposed System}
The proposed system is a standalone, web-based application focused on delivering a complete workflow from raw CSV data to a finalized BI dashboard. To automate data preparation for the dashboard, an intelligent sorting component using Gemini API is included. The system aims to demonstrate feasibility using realistic transaction data while remaining independent from direct banking or wallet API integrations. The core functionality is broken down as follows:

\begin{enumerate}
	\item \textbf{Custom Category Management:} The user will first define their own set of personal spending categories (e.g., Groceries, Transport, Health \& Fitness).
	\item \textbf{CSV Data Ingestion:} The user then uploads their transaction history as a standard CSV file. This serves as the primary data input mechanism.
	\item \textbf{Intelligent Sorting Component:} A smart classification component, utilizing Gemini API, processes the raw data. It analyzes transaction descriptions and human-like notes/comments, mapping each entry to the most relevant category from the user's own custom list. This automated step moves beyond simple rule-based mapping to understand context, preparing the data accurately for the BI layer.
	\item \textbf{Interactive BI Visualization:} Once sorted, the data is presented to the user through an interactive Business Intelligence (BI) dashboard. This dashboard is the primary output and uses Metabase~\cite{metabase2024bi} to visualize spending breakdowns tailored to the user's personal financial structure.
	\item \textbf{Validation with Public Data:} To demonstrate the effectiveness and real-world applicability of the sorting component, the system will be tested and validated using a relevant public dataset. This dataset will contain transaction-like entries with descriptions and comments, allowing for rigorous evaluation of the classification accuracy without relying on proprietary user data.
\end{enumerate}
By leveraging a sorting mechanism powered by Gemini API, adaptable to user categories and validated on realistic transaction data, the application delivers personalized and accurate insights through its core BI and visualization dashboard.

\section{Survey of the Related Applications}
To establish the necessity of this project, this survey analyzes two distinct application categories: the local digital wallets that create the data and the problem, and the global expense trackers that represent existing but flawed solutions.

\subsection{A. Market Survey — Local Digital Wallets (Direct Context)}
A review of the official documentation for Pakistan's leading digital wallets confirms that while they excel at payments, they generally lack a mature, integrated analytics tool.
\begin{itemize}
	\item \textbf{SadaPay, NayaPay, \& Easypaisa:} These platforms reliably provide users with the ability to download their account statements or view transaction histories, with some confirming CSV export availability~\cite{nayapay2023documents}. However, comparative reviews confirm that both SadaPay and NayaPay lack advanced in-app auto-categorization features~\cite{arshad2025sadapay}, and their feature sets do not include built-in BI or smart analytics layers. Indeed, product growth studies for Easypaisa have proposed introducing AI-powered expense tracking (such as Paisa Pulse) to incentivize premium subscriptions and solve this automated budgeting gap~\cite{nextleap2023easypaisa}.
	\item \textbf{JazzCash:} While also providing standard statements, JazzCash recently announced a Spending Tracker feature~\cite{jazzcash2023tracker}. Based on its public introduction, this appears to be a basic tracking tool for monitoring budgets. While a positive step, it is distinct from a full-scale BI layer and does not offer the intelligent, personalized categorization that is central to this proposal.
\end{itemize}
\textit{Implication.} The local market either offers no native analytics or, at best, provides basic tracking. This creates a clear analytics vacuum and validates the need for a separate application to provide the missing intelligence.

\subsection{B. Indirect Comparators — Global CSV-Import Tools}
Global platforms like YNAB and Mint demonstrate the feasibility of the CSV-import workflow but are poorly suited to the local context and fail to solve the core problem of personalization.
\begin{itemize}
	\item \textbf{High-Friction Workflow:} Tools like YNAB and Mint, while powerful, still require significant manual effort after a CSV import to clean up data and assign categories line-by-line. This does not solve the primary issue of laborious manual work.
	\item \textbf{Lack of Personalization and Intelligence:} These platforms are separate, often paid services that use rigid, generic categorization systems. They lack the intelligent sorting component needed to learn and adapt to a user's own custom financial categories, failing the key protein bar test case.
	\item \textbf{Misaligned Business Model:} Their primary model is direct bank API integration, which is not the standard in the Pakistani market. Their CSV import tools are a secondary feature, not the core product.
\end{itemize}
\textit{Implication.} The global tools prove the CSV-to-analytics model is viable but highlight a gap for a more intelligent, low-friction, and personalized solution tailored to the realities of the local market. This project is designed specifically to fill that gap.

\subsection{C. Comparative Analysis Framework}

\subsubsection{The Keyword Inflexibility Problem}
Traditional personal finance tracking software relies heavily on rigid string-matching routines. These rule-based parsers search for specific, exact keywords to match an entry to a category. While computationally fast, this framework fails when processing real-world transaction statements.

Local digital wallet statements frequently include irregular text strings, merchant branch codes, or variable transfer notes. A rule-based system looking for "KFC" will completely miss strings like "KFC-ISB-05" or "KFC Centaurus". This limitation forces users to constantly maintain complex mapping tables.

In addition to branch codes, transactions often append dynamic transaction numbers, payment dates, or terminal identifiers to the description field. These changing characters create unique strings for every purchase, even when the merchant is identical. Because a static keyword parser cannot match these variants, the system leaves many transactions uncategorized, requiring constant manual updates. Over time, the administrative overhead of adjusting rules makes these tracking tools too tedious for average users to maintain.

\subsubsection{Semantic Adaptation and Intent Recognition}
The Smart Transaction Sorter addresses this limitation by using natural language processing. Instead of treating text strings as exact keys, the system processes transaction descriptions as contextual data. It evaluates the vocabulary, merchant references, and user comments together.

The underlying language model recognizes the real-world semantic intent behind variations in text. This allows the system to map "Paid to INDRIVE ISLAMABAD" to a user-defined "Transport" category automatically. This semantic approach removes the need for constant manual rules, reducing setup effort for the end-user.

By looking at the complete text, the classifier captures contextual hints that rule-based systems ignore. For instance, payment descriptions containing words like transfer, sending, or receiving are evaluated alongside the payment type to determine the core intent of the transaction. This ensures that the system handles unusual or brand-new merchant names correctly without needing direct user intervention. It allows the classification system to adjust to new spending patterns automatically, making the initial setup much simpler.

\section{Modules \& Sub-modules}
The system architecture will be broken down into the following modules and sub-modules to ensure a clear separation of concerns and facilitate iterative development.

\subsection{User \& Category Management Module}
This module will handle all user-facing administrative tasks and personalization features.
\begin{itemize}
	\item \textbf{Authentication Submodule:} This submodule is responsible for secure user registration, login, logout, and session management to ensure data privacy.
	\item \textbf{Category CRUD Submodule:} This provides the interface and backend logic for users to Create, Read, Update, and Delete (CRUD) their personalized list of spending categories.
\end{itemize}

\subsection{Data Ingestion Module}
This module serves as the primary data entry point for transaction histories.
\begin{itemize}
	\item \textbf{File Upload Submodule:} This submodule provides the secure, user-facing web form for uploading transaction data in CSV format.
	\item \textbf{CSV Validation Submodule:} This is a backend service that parses the uploaded file and verifies its structural integrity, ensuring required columns (e.g., Date, Description, Amount) are present.
\end{itemize}

\subsection{Intelligent Sorting Module}
This module contains the core data processing and classification logic.
\begin{itemize}
	\item \textbf{Data Pre-processing Submodule:} This submodule cleans and standardizes the validated data, converting data types (e.g., string dates to date objects) to prepare it for classification.
	\item \textbf{Classification Engine Submodule:} This is the smart component that takes a transaction's text and the user's custom category list as input, and outputs the most probable category.
	\item \textbf{Database Persistence Submodule:} This submodule saves the processed transaction data, along with its newly assigned category, into the database.
\end{itemize}

\subsection{BI \& Visualization Module}
This module is responsible for presenting the final analytical output to the user.
\begin{itemize}
	\item \textbf{Token Authentication Submodule:} This is a backend service that compiles secure, signed JSON Web Tokens containing user session parameters using the HMAC-SHA256 algorithm.
	\item \textbf{Dashboard Rendering Submodule:} This is the frontend component that renders the securely embedded Metabase dashboard inside an iframe, applying the token parameters to display user-scoped charts and statistics.
\end{itemize}

\section{Primary Actors}
The system is designed for a single primary actor:
\begin{itemize}
	\item \textbf{The User:} An individual who wants to analyze their personal financial transaction history. The User will be able to perform the following actions:
	\begin{itemize}
		\item Register for an account and log in securely.
		\item Create and manage their own personalized list of spending categories.
		\item Upload their transaction history in CSV format.
		\item View the interactive BI dashboard with visualizations of their spending.
		\item (Optional) Review and manually correct any automated classifications.
	\end{itemize}
\end{itemize}

\section{Tools \& Techniques}
The project uses a modern, stable technology stack to ensure scalability and maintainability.
\begin{itemize}
	\item \textbf{Backend Framework:} Python with the Django Framework is used to build the server-side logic, manage URLs, and handle database interactions.
	\item \textbf{Database:} PostgreSQL is the confirmed database target for the final system because it supports reliable relational storage, indexing, and production deployment.
	\item \textbf{Data Processing:} Python CSV handling, Django forms, and model validation are used for efficient parsing, validation, and storage of uploaded transaction data.
	\item \textbf{Intelligent Sorting:} Gemini API is used to implement the classification component through structured prompts, JSON responses, confidence scoring, and category validation.
	\item \textbf{Frontend Visualization:} Metabase is used as the Business Intelligence layer and is embedded into the Django interface through a secure dashboard URL.
	\item \textbf{Web Technologies:} The user interface is built with Django templates, HTML, CSS, JavaScript, Bootstrap, and Jazzmin for the administrative interface.
\end{itemize}

\section{Process Model \& Design Approach}
The project will be developed using an Iterative and Incremental process model. This approach is well-suited for a modular system, as it allows for the development and testing of each component in successive cycles or iterations.

The development will proceed by building and testing each module individually before integrating them into a cohesive whole. This allows for early feedback, easier debugging, and a more manageable workflow. The design approach will follow the Model-View-Template (MVT) architectural pattern, which is native to the Django framework and provides a clear separation between data logic, business logic, and user interface.

\section{Project Plan}
The project timeline is scheduled over approximately seven months (30 weeks), visualized in the Gantt chart in Figure \ref{fig:gantt} below. The plan prioritizes documentation and core feature development first, followed by the integration of the intelligent sorting component and finalization.

\begin{figure}[H]
	\centering
	\resizebox{\textwidth}{!}{
		\begin{ganttchart}[
			vgrid,
			hgrid,
			bar/.append style={fill=blue!50},
			milestone/.append style={fill=red!70}
			]{1}{30}
			\gantttitle{Project Timeline (Weeks)}{30} \\
			\gantttitlelist{1,...,30}{1} \\
			
			% Phase 1: Documentation & Initial Setup
			\ganttgroup{Phase 1: Planning \& Design}{1}{8} \\
			\ganttmilestone[name=p1]{Proposal Submission}{1} \ganttnewline
			\ganttlinkedbar[name=srs]{SRS Documentation}{2}{3} \ganttnewline
			\ganttlinkedbar[name=uc]{Use Case Diagrams}{3}{5} \ganttnewline
			\ganttlinkedbar[name=ui]{Initial UI/UX Design}{5}{8} \ganttnewline
			
			% Phase 2: Core CSV & BI Development
			\ganttgroup{Phase 2: Core Feature Development}{9}{18} \\
			\ganttlinkedbar[name=data]{Develop Data Ingestion (CSV)}{9}{12} \ganttnewline
			\ganttlinkedbar[name=db]{Database \& Backend Setup}{11}{14} \ganttnewline
			\ganttlinkedbar[name=bi]{Develop BI \& Visualization Module}{13}{18} \ganttnewline
			
			% Phase 3: AI Integration & Testing
			\ganttgroup{Phase 3: AI Integration}{19}{25} \\
			\ganttlinkedbar[name=ai]{Develop Intelligent Sorting Module}{19}{22} \ganttnewline
			\ganttlinkedbar[name=integ]{Integrate Sorting with Main App}{23}{25} \ganttnewline
			
			% Phase 4: Finalization & Reporting
			\ganttgroup{Phase 4: Finalization \& Reporting}{26}{30} \\
			\ganttlinkedbar[name=test]{End-to-End Testing \& Bug Fixing}{26}{28} \ganttnewline
			\ganttlinkedbar[name=doc]{Final Report Documentation}{28}{29} \ganttnewline
			\ganttmilestone[name=p2]{Final Presentation}{30}
			
			% Links
			\ganttlink{p1}{srs}
			\ganttlink{srs}{uc}
			\ganttlink{uc}{ui}
			\ganttlink{ui}{data}
			\ganttlink{data}{db}
			\ganttlink{db}{bi}
			\ganttlink{bi}{ai}
			\ganttlink{ai}{integ}
			\ganttlink{integ}{test}
			\ganttlink{test}{doc}
			\ganttlink{doc}{p2}
		\end{ganttchart}
	}
	\caption{Project development timeline presented as a Gantt chart.}
	\label{fig:gantt}
\end{figure}

\section{Socio-Economic Context of FinTech in Pakistan}

\subsection{The Shift to Mobile Banking Ecosystems}
The financial landscape in Pakistan has changed rapidly due to the growth of branchless banking platforms. Services like Easypaisa, JazzCash, NayaPay, and SadaPay have enabled millions of individuals to manage money digitally. This has led to a massive increase in transaction volumes across a variety of user demographics.

Despite this growth in digital transactions, financial literacy tools have not kept pace. Most wallet providers optimize their interfaces for quick payments rather than long-term asset management. Users are left with unstructured text logs that do not offer clear visibility into their personal cash flows.

The absence of native tracking features makes it hard for users to build positive saving habits. Digital statements show when and where money was spent, but they don't help users understand their spending patterns. This gap in service creates a clear need for secondary analytical tools that can organize cash flows.

\subsection{The Cognitive Load of Unstructured Financial Data}
Raw data statements exported as CSV files create an analytical barrier for everyday consumers. Disorganized rows of dates and numeric strings make manual categorization tedious and frustrating. This data complexity often prevents users from maintaining consistent personal budgets.

The lack of a domestic personal finance management standard highlights the need for an automated layer. A standalone utility that cleans irregular formats and presents data visually makes financial tracking accessible. This system provides a practical blueprint for improving consumer financial awareness in the local market.

By automating the sorting task, the system removes the mental effort required to manage personal accounts. Users no longer need to spend hours grouping transactions by hand or deciphering vague merchant text. This convenience encourages regular budget reviews, helping individuals make better financial decisions.

\subsection{Empirical Challenges in Local Financial Text Processing}
The development of automated financial tracking utilities within the Pakistani fintech ecosystem faces significant text-processing challenges. Unlike international banking data, which often adheres to standardized merchant category codes and clean identification logs, local transaction logs from mobile wallets and domestic bank statements are notoriously unpredictable. These records frequently combine multiple languages, structural variations, and random numbers within a single description string.

The core text processing issues can be broken down into three distinct semantic categories:
\begin{itemize}
    \item \textbf{Alphanumeric Fragmentation and Dynamic IDs}: Local statements regularly interweave merchant identifiers with variable transactional sequences. For example, a single merchant transaction might append unique billing serials, date numbers, or terminal codes directly to the shop name. Because these codes change with every single purchase, traditional string-matching utilities fail completely. The text parser cannot determine where the stable merchant identifier ends and where the dynamic transaction counter begins, creating severe classification friction.
    \item \textbf{Bilingual Inconsistencies and Romanized Script}: Financial text logs in the local market rely heavily on Romanized Urdu phrasing mixed with standard English financial terms. Transaction notes routinely feature terms like bhai, shukriya, khana, or udhaar right alongside terms like funds transfer or bill payment. Furthermore, phonetic variations introduce massive structural differences. The same target concept or location can be spelled in multiple ways across different applications, requiring an analytics system that can recognize underlying semantic intent rather than relying on strict spelling matches.
    \item \textbf{Cryptic Merchant Abbreviations}: Point of sale machines and digital payment portals frequently truncate merchant listings into short, compressed character chains to fit narrow statement interfaces. A clear storefront name is often reduced to a confusing set of initials combined with local city or branch indicators. Without an intelligent natural language processing layer that evaluates the complete contextual signature of the text string, basic keyword rules treat these shortened tags as completely separate entities, leaving them entirely uncategorized.
\end{itemize}
